\documentclass[11pt,a4paper]{article}
\usepackage[margin=26mm]{geometry}
\usepackage[T1]{fontenc}
\usepackage{lmodern,microtype,amsmath,amssymb,amsthm,mathtools,booktabs,array}
\usepackage[hidelinks]{hyperref}
\newcommand{\Qset}{\mathcal Q_{222}}
\newcommand{\R}{\mathbb R}
\newcommand{\Id}{\mathbb I}
\newcommand{\conv}{\operatorname{conv}}
\newcommand{\Sym}{\operatorname{Sym}}
\newcommand{\tr}{\operatorname{tr}}
\newcommand{\res}{\operatorname{res}}
\newcommand{\sos}{\Sigma T^2}
\newtheorem{theorem}{Theorem}
\newtheorem{proposition}[theorem]{Proposition}
\newtheorem{lemma}[theorem]{Lemma}
\newtheorem{corollary}[theorem]{Corollary}
\theoremstyle{definition}
\newtheorem{definition}[theorem]{Definition}
\setlength{\emergencystretch}{2em}
\allowdisplaybreaks[1]
\title{Finite semidefinite representation of the simplest quantum correlation set}
\author{}
\date{}
\hypersetup{pdftitle={Finite semidefinite representation of the simplest quantum correlation set}}
\begin{document}
\makeatletter
\begin{center}
{\LARGE\@title\par}
\end{center}
\makeatother
\vspace{1em}
\begin{abstract}
The simplest Bell experiment allows two binary measurements on each of two separated systems. We prove that its complete quantum correlation set, including all marginal expectations and joint correlators, admits a finite semidefinite representation. Consequently, one fixed semidefinite program suffices to optimize every Bell functional in this scenario exactly. The proof reduces arbitrary-dimensional quantum realizations to real two-qubit correlations and constructs exact sums of squares through the singular limits in which measurements become compatible. The inequality coefficients and measurement parameters vary independently throughout the argument. A real-algebraic criterion converts these certificates into a finite global representation, whose matrices can be chosen with real algebraic entries.

\end{abstract}

\section{Introduction}

Which joint outcome probabilities can two separated quantum systems produce?
Even when each party chooses between two binary measurements, the answer
depends on both the correlations between the outcomes and their individual
biases. Together these quantities form an eight-dimensional convex set.
Its supporting hyperplanes give all quantum bounds on Bell inequalities in
this experiment.

Semidefinite programming provides a systematic way to bound these quantum
values. The Navascu\'es--Pironio--Ac\'in hierarchy imposes positivity on
successively larger moment matrices~\cite{NPA}. In the two-input,
two-output scenario, a different finite description follows from the
reduction to two-qubit experiments~\cite{Masanes}: the quantum set is the
convex hull of a compact polynomial image. Its extreme points also admit
an analytical characterization~\cite{BarizienBancal}. These descriptions
make the set accessible to real algebraic geometry. They do not, by
themselves, provide a finite linear matrix inequality for membership.
Indeed, compact convex semialgebraic sets can fail to admit any
semidefinite lift~\cite{Shadows}.

We prove that the complete quantum set in this scenario does admit such
a lift. There is one finite positive-semidefinite constraint, with
auxiliary real variables, whose projection is exactly the set of all
eight expectations. The same feasible region therefore gives the exact
quantum value of every Bell functional. The matrices can be chosen with
real algebraic entries.

The difficulty occurs near compatible measurements. An inequality may
approach a singular supporting hyperplane while an independent quantum
experiment approaches a product-state or commuting-measurement boundary.
A certificate valid at a fixed interior point need not remain valid
through these simultaneous limits. We keep the coefficients of the
inequality and the parameters of the experiment in separate real closed
fields. Exact sums of squares in their tensor product then express the
positivity needed for a finite lift. Scheiderer's tensor-evaluation
criterion~\cite{Scheiderer} turns this universal positivity statement
into semidefinite representability.

We develop the physical reduction, the tensor criterion and the global
assembly below. The local proof has two parts: regular contacts are
controlled by curvature, whereas the compatible-measurement limits
require explicit rescalings of the complete state polynomial. The
accompanying \emph{Complete proofs} supplies every local calculation
and the source partitions used in the assembly.

\section{The quantum set and the main theorem}

Alice chooses $x\in\{0,1\}$ and Bob chooses $y\in\{0,1\}$.
Their outcomes $a,b$ take values in $\{+1,-1\}$. A quantum behavior has
the form
\begin{equation}
 p(a,b\mid x,y)=\tr\bigl[\rho(M_{a|x}\otimes N_{b|y})\bigr],
 \label{eq:behavior}
\end{equation}
where $\rho$ is a bipartite density operator and each measurement is a
binary positive-operator-valued measurement. Arbitrary finite local
Hilbert-space dimensions are allowed. Define the observables
$A_x=M_{+|x}-M_{-|x}$ and $B_y=N_{+|y}-N_{-|y}$. Nonsignalling makes
the four marginal expectations independent of the other party's input.
We use the eight coordinates
\begin{equation}
 q=(a_0,a_1,b_0,b_1,E_{00},E_{01},E_{10},E_{11}),\qquad
 a_x=\langle A_x\rangle,\quad b_y=\langle B_y\rangle,
 \quad E_{xy}=\langle A_xB_y\rangle.
 \label{eq:coordinates}
\end{equation}
They recover every probability through
\begin{equation}
 p(a,b\mid x,y)=\frac{1+a a_x+b b_y+abE_{xy}}4.
 \label{eq:probabilities}
\end{equation}
Let $\Qset\subset\R^8$ denote the set obtained from
Eq.~\eqref{eq:behavior}.

\begin{theorem}[Finite semidefinite representation]
\label{thm:main}
There exist finite integers $N,p$ and real symmetric matrices
$L_0,L_1,\ldots,L_8,M_1,\ldots,M_p$ of order $N$ such that
\begin{equation}
 \Qset=\left\{q\in\R^8:\ \exists y\in\R^p,\quad
 L_0+\sum_{\nu=1}^{8}q_\nu L_\nu+
                   \sum_{j=1}^{p}y_jM_j\succeq0\right\}.
 \label{eq:lift}
\end{equation}
All entries of these matrices may be chosen to be real algebraic numbers.
\end{theorem}

Thus $\Qset$ is a spectrahedral shadow. For any Bell functional
$\beta(q)=\beta_0+\sum_\nu\beta_\nu q_\nu$, its maximum over
Eq.~\eqref{eq:lift} is exactly its quantum maximum. Compactness of
$\Qset$, proved next, ensures that this maximum is attained.
The theorem establishes existence of a finite representation; evaluating
its matrices and finding a useful matrix size are subsequent constructive
questions.

\section{A finite polynomial description of all quantum behaviors}
\label{sec:physical}

The field-extension argument requires a fixed finite description of
$\Qset$. We obtain one by reducing the measurements before taking the
convex hull.

For an effect $0\preceq E\preceq I$, the matrix
\begin{equation}
 \widehat P_E=\begin{pmatrix}
 E&\sqrt{E(I-E)}\\ \sqrt{E(I-E)}&I-E
 \end{pmatrix},\qquad Vh=(h,0),
 \label{eq:dilation}
\end{equation}
is a projection and satisfies $V^*\widehat P_EV=E$. The embedding $V$
is the same for both settings of a party. Applying it on both sides
preserves all eight expectations. We may therefore use observables
whose squares are the identity.

Two such observables have a common decomposition into invariant blocks
of dimension at most two. To see this, let $U,V$ be Hermitian involutions.
The unitary $C=UV$ satisfies $UCU=C^*$. An eigenvector of $C$ with
nonreal eigenvalue $\lambda$ pairs with its image under $U$, which has
eigenvalue $\overline\lambda$; their span is invariant under both
observables. On the $\lambda=\pm1$ eigenspaces, $V=\pm U$, and
one-dimensional blocks suffice. Repeating this decomposition gives the
claimed blocks.

Within a two-dimensional block a choice of basis and phase puts the
measurements in the form
\begin{equation}
 A_0=B_0=Z,\qquad A_1=c_AZ+s_AX,\qquad B_1=c_BZ+s_BX,
 \quad c_A^2+s_A^2=c_B^2+s_B^2=1,
 \label{eq:qubit}
\end{equation}
where $X,Z$ are the Pauli matrices. One-dimensional blocks embed in this
description at $s_A=0$ or $s_B=0$. Since the measurements are block
diagonal, their expectations are convex combinations of two-qubit
expectations. The matrices in Eq.~\eqref{eq:qubit} are real. For each
real symmetric measurement matrix $M$,
\begin{equation}
 (r+is)^*M(r+is)=r^TMr+s^TMs.
 \label{eq:real-state}
\end{equation}
Normalizing the nonzero real and imaginary parts separately therefore
reduces complex states to mixtures of real pure states.

Let $\Phi:S^1\times S^1\times S^3\to\R^8$ be the expectation map
defined by Eq.~\eqref{eq:qubit} and a real unit state
$\psi\in\R^4$. Every coordinate of $\Phi$ is polynomial.
The preceding reduction and the physical realization of a finite mixture
by direct sums give
\begin{equation}
 \Qset=\conv\Phi(S^1\times S^1\times S^3)
 =\left\{\sum_{i=1}^{9}w_i\Phi(z_i):
 w_i\ge0,\ \sum_iw_i=1,\ z_i\in S^1\times S^1\times S^3\right\}.
 \label{eq:nine}
\end{equation}
The bound of nine is the Carath\'eodory bound in eight dimensions.
It limits the number of atoms in a representation; the states and
measurement angles still range continuously.

Equation~\eqref{eq:nine} is a polynomial formula with rational
coefficients and finitely many variables. Its domain is compact over
$\R$, so $\Qset$ is compact. Projection of a semialgebraic set is
semialgebraic~\cite{Basu}. We can consequently transfer this same finite
formula to any real closed extension $K$ of $\R$; the resulting set is
denoted $(\Qset)_K$. This supplies the precise meaning of a quantum
point over $K$ in the next section.

\section{From Bell bounds to a finite lift}
\label{sec:tensor}

Write $P$ for the cone of affine functions nonnegative on $\Qset$.
A Bell bound $\beta(q)\le b$ gives such a function,
$f(q)=b-\beta(q)$. We call $f$ a Bell slack. Over a real closed
extension $K$, the transferred cone $P_K$ consists exactly of the
affine functions nonnegative on $(\Qset)_K$. This follows by
transferring the finite formula
$\forall q\,[q\in\Qset\Rightarrow f(q)\ge0]$.

To keep the inequality and the experiment independent, choose two
labelled real closed extensions $K_c,K_t$ of $\R$ and put
$T=K_c\otimes_{\R}K_t$. For $f=f_0+\sum_\nu f_\nu q_\nu$
with coefficients in $K_c$, and $q\in(\Qset)_{K_t}$, define
\begin{equation}
 [f,q]=f_0\otimes1+\sum_{\nu=1}^{8}f_\nu\otimes q_\nu.
 \label{eq:tensor}
\end{equation}
An element of $T$ is a finite sum of separated products. We seek an
identity $[f,q]=\sum_k h_k^2$ in this ring. The fields may be labelled
copies of the same field, but their elements retain their separate
positions in the tensor product.

\begin{theorem}[Universal sum of squares]
\label{thm:tensor}
For every $K_c,K_t\supset\R$, every $f\in P_{K_c}$ and every
$q\in(\Qset)_{K_t}$,
\begin{equation}
                         [f,q]\in\Sigma T^2.
 \label{eq:universalsos}
\end{equation}
Here $\Sigma T^2$ is the set of finite sums of squares in $T$.
\end{theorem}

Scheiderer's Corollary~3.19~\cite{Scheiderer} states that a closed
convex semialgebraic set $S\subset\R^n$ has a finite semidefinite
lift precisely when the corresponding tensor evaluations are sums of
squares for every real closed extension $K/\R$, every valid affine
function over $K$, and every point of $S_K$. Applying
Theorem~\ref{thm:tensor} to two labelled copies of each $K$ gives this
condition for $\Qset$. Together with Eq.~\eqref{eq:nine}, it proves
the real-coefficient part of Theorem~\ref{thm:main}.

The criterion permits the number and tensor ranks of the squares to
depend on the point and the affine function. Its conclusion supplies
a common finite lift. It is therefore enough to construct exact
certificates throughout the source and testing domains. The next
two sections make those domains and their certificates explicit.

\section{Two families of supporting inequalities}
\label{sec:sources}

We first reduce the cone of Bell slacks to a small collection of
coefficient patterns. Choose the CHSH expression
$S=E_{00}-E_{01}+E_{10}+E_{11}$ and set $H=1-S/2$.
At the corresponding Popescu--Rohrlich vertex, $H=-1$ and the eight
losing-event probabilities vanish. Order these events as
\begingroup\setlength{\arraycolsep}{3pt}
\begin{equation}
 \begin{array}{c|cccccccc}
 \text{weight}&\lambda&n&m&C&A&D&\alpha&\beta\\
 \text{event}&(++|01)&(+-|00)&(+-|11)&(-+|10)&(-+|00)&(-+|11)&(--|01)&(+-|10)
 \end{array}
 \label{eq:events}
\end{equation}
\endgroup
The letters on the left denote nonnegative weights, not probabilities.
Every valid slack negative at that vertex has the unique form
\begin{equation}
 f=dH+4\sum_i c_i p_i,\qquad d>0,\quad c_i\ge0.
 \label{eq:source}
\end{equation}
The coefficients follow by evaluating $f$ at that vertex and its
eight deterministic CHSH-facet neighbors. If $f$ is nonnegative at
every PR vertex, it is nonnegative on the nonsignalling polytope and
is a nonnegative combination of event probabilities. Those slacks
already have separated square certificates.

For the remaining slacks, introduce half-angles $u,v$ and a real
homogeneous state $(W,X,Y,Z)$. The measurement circle coordinates are
$(1-u^2,2u)/(1+u^2)$ and $(1-v^2,2v)/(1+v^2)$.
Multiplication by the positive measurement denominators and state norm
gives the complete polynomial
\begin{align}
 F={}&4\lambda D_u(W+vX)^2+4nD_uD_vX^2
          +4m(vW-X+uvY-uZ)^2\notag\\
 &+4CD_v(uW-Y)^2+4AD_uD_vY^2
          +(4D+2d)(uW+uvX-Y-vZ)^2\notag\\
 &+(4\alpha+2d)D_u(vY-Z)^2
          +(4\beta+2d)D_v(X+uZ)^2-2dD_uD_vZ^2,
 \label{eq:parent}
\end{align}
where $D_u=1+u^2$ and $D_v=1+v^2$.
Thus $F=D_uD_v\|\psi\|^2 f(\Phi)$ in one field, and the same
identity holds with separated coefficients in $T$. Homogeneous angle
pairs give the polynomial at the missing projective endpoints.

The negative term in Eq.~\eqref{eq:parent} is the only obstruction
to an immediate event-square decomposition. Let $P(u,v)$ be its
four-by-four state matrix. For $d>0$ and $uv\ne0$, the restriction
to $Z=0$ contains the positive definite form
\begin{equation}
 2d\bigl[(uW+uvX-Y)^2+D_uv^2Y^2+D_vX^2\bigr].
 \label{eq:stateblock}
\end{equation}
Hence $P\succeq0$ is equivalent to $\det P\ge0$. Its determinant
has a useful structure:
\begin{equation}
 \det P=16D_u^2D_v^2u^2v^2 L(u^2,v^2),\qquad
 L(U,V)=h+\sum_{(i,j)\ne(1,1)}q_{ij}U^{i-1}V^{j-1},
 \quad q_{ij}\ge0,
 \label{eq:determinant}
\end{equation}
where $(i,j)\in\{0,1,2\}^2\setminus\{(1,1)\}$.
The eight coefficients are displayed in Complete proofs, Chapter~4,
Eq.~(SG6). Their nonnegativity makes $L(e^s,e^t)$ convex.
Consequently, a stationary zero at finite positive angles is a global
minimum. Together with Eq.~\eqref{eq:stateblock}, this proves that
the six contact equations
\begin{equation}
 P(u_c,v_c)\psi_c=0,\qquad
 \psi_c^TP_u(u_c,v_c)\psi_c=
 \psi_c^TP_v(u_c,v_c)\psi_c=0
 \label{eq:contact}
\end{equation}
already certify physicality of a nonnegative coefficient tuple.

Include the event-face boundary $d=0$ and normalize all nine
nonnegative coefficients in Eq.~\eqref{eq:source} to sum to one.
The physical tuples form a compact convex body. At an extreme point
with $d>0$ whose Laurent
polynomial is coercive, the minimum of $L$ must be zero; otherwise
small opposite coefficient perturbations would remain physical.
At this contact, Eq.~\eqref{eq:contact} gives six linear equations
in the source coefficients. Seven positive event weights would leave
a nonzero perturbation preserving those equations and the normalization,
again contradicting extremality. The compatibility of the omitted
events then leaves exactly two patterns, up to Bell relabellings:
\begin{equation}
\begin{array}{ll}
\mathsf C_1:&\lambda=n=0,\quad d,m,C,A,\alpha>0,
                       \quad D,\beta\ge0,\quad D+\beta>0,\\[2pt]
\mathsf C_2:&\lambda=\beta=0,\quad d,m,C,A,D,\alpha>0,
                       \quad n\ge0.
\end{array}
\label{eq:families}
\end{equation}
The noncoercive supports reduce to two other patterns,
$\lambda=m=0$ or $\lambda=\alpha=0$. Their source cones have explicit
matrix lifts: a six-by-six Gram matrix with three auxiliaries in the
first case, and two three-by-three blocks with one auxiliary in the
second. Complete proofs, Chapter~4, derives both lifts and their
attained dual representations. Pairing a source lift with its testing
dual gives separated squares directly.

Every normalized source is a mixture of at most nine extreme sources.
This finite decomposition, the contact equations and the support
classification all transfer to the coefficient field. It remains to
certify the full polynomial~\eqref{eq:parent} for the two families
in Eq.~\eqref{eq:families} at arbitrary independent tests.

\section{Exact certificates at regular and singular contacts}
\label{sec:local}

\subsection{A positive real Gram matrix controls infinitesimal errors}

For $K_i\supset\R$, let $B_i$ be its ring of bounded elements and
$\mathfrak m_i$ its ideal of infinitesimals. Every bounded element has
a unique real standard part. This defines a residue map
$\rho:B_c\otimes B_t\to\R$ with kernel
$\mathfrak m_c\otimes B_t+B_c\otimes\mathfrak m_t$.
Thus a zero-residue coefficient is a finite sum of separated terms,
each with an infinitesimal factor.

The following elementary identity turns a positive residue into exact
squares. If $x\in K_c$ is infinitesimal, $y\in K_t$ is bounded,
and $\epsilon,M$ are positive reals with $M>|y|$, then
\begin{equation}
 \epsilon+x\otimes y=
 \tfrac12\bigl[(\epsilon/M+x)\otimes(M+y)
              +(\epsilon/M-x)\otimes(M-y)\bigr].
 \label{eq:reserve}
\end{equation}
Each pure-field factor on the right is positive, so each product is
an elementary tensor square. A positive real residue can be split
among finitely many such terms.

\begin{lemma}[Bounded Gram matrix]
\label{lem:gram}
If $G\in\Sym_r(B_c\otimes B_t)$ and $\rho(G)\succ0$, then
$z^TGz\in\Sigma T^2$ for every $z\in T^r$.
\end{lemma}
\begin{proof}
A fixed real congruence reduces the residue to the identity. Choose
$\delta>0$ with $(r-1)\delta<1$. The exact decomposition
\begin{align}
 G={}&\sum_i\bigl(G_{ii}-(r-1)\delta\bigr)e_ie_i^T\notag\\
 &+\sum_{i<j}\frac{\delta+G_{ij}}2(e_i+e_j)(e_i+e_j)^T
 +\sum_{i<j}\frac{\delta-G_{ij}}2(e_i-e_j)(e_i-e_j)^T
 \label{eq:gramframe}
\end{align}
has coefficients of positive residue. Apply Eq.~\eqref{eq:reserve}
to each coefficient and evaluate on $z$. The roots in $z$ may be
unbounded; only the Gram coefficients require boundedness.
\end{proof}

\subsection{Contacts with noncommuting measurement residues}

Suppose both local measurement pairs remain noncommuting after taking
real residues. At a zero of the Bell slack the residue state must be
entangled, since a product state gives a strict CHSH bound at these
angles. An entangled two-qubit state is an invertible bilinear form
on the local rays. Each local ray can therefore be orthogonal to at
most one remote ray. The vanishing losing events form a matching
in the cycle of eight events, so at most four of them can vanish.

This finite constraint has an important consequence for unequal source
scales. Retain the event amplitudes with the largest coefficients as
exact state coordinates. If the next source scale has a positive
constant term, Lemma~\ref{lem:gram} applies. If the signed coefficient
is visible at that scale, the constrained contact has a positive
Hessian. The angular part follows from Eq.~\eqref{eq:determinant}:
at a stationary zero,
\begin{equation}
 \nabla^2_{\log U,\log V} L
       =\sum_{e} q_e U^{e_1}V^{e_2}\,ee^T.
 \label{eq:curvature}
\end{equation}
The active exponent vectors span the plane. The event supports that
would reduce their span have an explicit negative CHSH-core direction
and are excluded by physicality. The positive transverse state block
then gives the full constrained Hessian.

To lift this curvature through independent infinitesimals, write the
complete local polynomial as
\begin{equation}
 F(x,y)=\sum_i a_i b_i(y)x_i^2+\sigma G(x,y),\qquad
 \varepsilon_i=\sigma/a_i>0\text{ infinitesimal}.
 \label{eq:weighted}
\end{equation}
The ratios $\varepsilon_i$ may differ arbitrarily. Substituting
$x_i=\varepsilon_i s_i$ into the stationary equations gives a
nonsingular real residue Jacobian. The semialgebraic implicit-function
theorem supplies an exact center $(x_c,y_c)$ in the coefficient field.
Finite Taylor expansion there gives
\begin{equation}
 F/\sigma=a+h^TQh,\qquad a\ge0,\quad\rho(Q)\succ0,
 \qquad h=\bigl((x_i-x_{ci})/\sqrt{\varepsilon_i}\bigr)_i
                         \mathbin{\oplus}(y-y_c).
 \label{eq:weightedgram}
\end{equation}
The vanishing first derivatives remove linear terms. Every higher
coefficient acquires a bounded or infinitesimal pure-field factor
after the displayed rescaling. Lemma~\ref{lem:gram} gives exact
squares of the entire polynomial.

If the signed coefficient remains infinitesimal at the next scale,
another positive event amplitude must vanish. Add it to the retained
constraints. Each repetition adds an edge to the matching, so the
process stops after at most four edges. The four-edge case has an
explicit weighted completion and a strictly positive remaining
tangential curvature. Complete proofs, Chapter~6, proves this
terminal case for arbitrary ratios of the four coefficients.

\subsection{Limits of compatible measurements}

When a measurement pair becomes compatible in residue, the curvature
in Eq.~\eqref{eq:curvature} may vanish. The exact contact equations
then determine the rescaling needed to keep the lost directions.
We first solve those equations for the source coefficients and prove
the strict signs of their denominators. The derivation uses a connected
graph of real contacts, explicit nonzero minors and event-deletion
arguments. It also includes the permitted zero coefficients in
Eq.~\eqref{eq:families}; see Complete proofs, Chapter~5.

These coordinates divide the source domain into finitely many cases.
The first division uses the signed coefficient relative to the sum of
all coefficients. A positive residue gives the ordinary edge and
corner cases. An infinitesimal residue gives the two families
$\mathsf C_1,\mathsf C_2$ with further exact source inequalities.
Every positive quantity is either infinitesimal, a bounded unit or
unbounded in its own field. Literal zero coefficients are kept as
separate boundaries. This classification exhausts the source domain
without comparing the rates in the two fields.

For each case, choose state deviations $\xi$ from the relevant contact
graph and retain the identity
\begin{equation}
 F=F_0+2\sum_i\xi_iF_i+\sum_{i,j}\xi_i\xi_jF_{ij}.
 \label{eq:fullnormal}
\end{equation}
The scalar $F_0$, all first-normal coefficients $F_i$ and the whole
quadratic state form are used together. Contact identities express
the vanishing parts through finitely many polynomial roots. Monic
division preserves the required source factors; it never divides by
a difference between source and testing parameters. Rescaling these
roots gives bounded Gram coefficients with a strict real residue.
Complementary state charts and measurement directions are treated
directly from Eq.~\eqref{eq:parent}.

One identity illustrates how the source factors survive this division.
Let $k$ be a positive source infinitesimal and $s$ a positive source
unit. For independent positive testing coordinates $u,v$, put
\begin{equation}
 U_1=u^2v,\qquad U_2=kuv^2/s,\qquad U_3=k^2s,\qquad g=kuv.
 \label{eq:circuitroots}
\end{equation}
Then $U_1U_2U_3=g^3$, and the exact identities
\begin{equation}
 g(u-k)=k(U_1-g),\qquad g(v-s)=s(U_2-g)
 \label{eq:conductor}
\end{equation}
replace contact differences by the roots already present in the
three-monomial sum-of-squares circuit. The first replacement supplies
the infinitesimal factor $k$. In the second, $s$ is a unit, so the
corresponding Taylor coefficients must themselves be infinitesimal;
the local source formulas establish this before the identity is used.
The first contact derivatives cancel the remaining linear terms.
Neither $u-k$ nor $v-s$ is inverted. Complete proofs, Appendix~B,
gives the resulting finite Gram correction and the full state-normal
extension.

The resulting statement is the local input needed for the global proof.

\begin{proposition}[Complete canonical certificate]
\label{prop:canonical}
Let $K_c,K_t\supset\R$ be labelled real closed fields. Let the
coefficients of Eq.~\eqref{eq:parent} belong to either family in
Eq.~\eqref{eq:families} and admit a contact satisfying
Eq.~\eqref{eq:contact} with $u_c,v_c>0$ and $\psi_c\ne0$.
Then its full homogeneous tensor
evaluation is a finite sum of squares for every measurement pair
and every state in $K_t$.
\end{proposition}

Chapter~7 of Complete proofs assembles the local certificates.
Appendix~A treats bounded source coefficients; Appendices~B and C
establish the two remaining contact families, including their
complementary testing charts. Several singular cases share one
quadratic completion. In $\mathsf C_1$, one contact-division argument
allows the source and testing coordinates to have either zero or
nonzero residue. Its coefficients retain an infinitesimal factor
after every division. In $\mathsf C_2$, the four vanishing event
gaps enter through one radial scale; an exact current relation
cancels the cross term before that scale is divided out. The
remaining source domains have the explicit positive matrix
completions given there. Source and testing perturbations remain
independent throughout.

\section{Global assembly and algebraic coefficients}
\label{sec:assembly}

We can now prove Theorem~\ref{thm:tensor}. Decompose a valid affine
function into finitely many event, closed-face and canonical contact
sources. Event sources are separated squares. The explicit closed-face
lifts give positive-semidefinite matrices $X_c$ and $Y_t$ over their
respective fields. Their pairing is a sum of squares because
\begin{equation}
 \sum_{i,j}(X_c)_{ij}\otimes(Y_t)_{ij}
   =\sum_{a,b}\left(\sum_i U_{ia}\otimes V_{ib}\right)^2,
 \qquad X_c=UU^T,\quad Y_t=VV^T.
 \label{eq:pairing}
\end{equation}
The factorizations exist over real closed fields, including for
singular matrices. Proposition~\ref{prop:canonical} handles the
remaining source families. Dividing out the measurement denominators
and the state norm multiplies each root only by positive inverse
square roots in the testing field. Source normalizations are restored
in the coefficient field.

For completeness, signs and endpoints do not introduce further cases.
Local changes of basis give positive finite half-angles. If one pair
commutes literally, the CHSH operator is $\pm2$ times a reflection
on each local block, so $H\succeq0$ in the testing field. At positive
half-angles the CHSH core is also nonnegative when $WZ\le0$, by its
explicit two-state Schur form. The remaining states have $W/Z>0$
and are covered by Proposition~\ref{prop:canonical}. Homogeneous
angle coordinates retain both projective endpoints.

Write $f=\sum_j a_jf_j$ with $a_j\ge0$ in $K_c$, and use
Eq.~\eqref{eq:nine} to write $q=\sum_iw_iq_i$ in $K_t$.
For each source and testing atom, let $[f_j,q_i]=\sum_kh_{ijk}^2$.
Then
\begin{equation}
 [f,q]=\sum_{i,j,k}
           \bigl[(\sqrt{a_j}\otimes\sqrt{w_i})h_{ijk}\bigr]^2.
 \label{eq:assembly}
\end{equation}
The testing weights sum to one and therefore preserve the affine
constant. Zero weights give zero roots. This proves
Theorem~\ref{thm:tensor}, and the tensor criterion proves
Eq.~\eqref{eq:lift}.

It remains to justify the algebraic entries asserted in
Theorem~\ref{thm:main}. Let $\R_{\mathrm{alg}}$ be the real algebraic
numbers. This is a real closed subfield of $\R$.
Fix the finite sizes $N,p$ supplied by the existence theorem.
Equality between Eq.~\eqref{eq:nine} and a projected pencil of these
sizes is a first-order sentence over $\mathbb Q$:
\begin{equation}
 \exists L_0,\ldots,L_8,M_1,\ldots,M_p\quad
 \forall q\quad
 \left[q\in\Qset\ \Longleftrightarrow\
 \exists y\;L_0+\sum_\nu q_\nu L_\nu+\sum_jy_jM_j\succeq0\right].
 \label{eq:algebraic}
\end{equation}
Here membership uses the finite polynomial formula~\eqref{eq:nine},
and matrix positivity uses finitely many principal minors. The sentence
is true over $\R$ and hence over $\R_{\mathrm{alg}}$. Choose its
matrix coefficients there. Elementary transfer back to $\R$ gives
Eq.~\eqref{eq:lift} for every real $q$, with those same algebraic
entries. This completes the proof of Theorem~\ref{thm:main}.

\section{An explicit projection}

A two-dimensional projection illustrates what an evaluated representation
looks like. Set $t=a_0$ and $s=E_{00}+E_{01}+E_{10}-E_{11}$.
The exact projection is
\begin{equation}
 \mathcal T=\{(t,s):|t|\le1,\quad t^2+s^2/4\le2\}.
 \label{eq:tiltedset}
\end{equation}
It has the order-four linear matrix inequality
\begin{equation}
 \begin{pmatrix}\sqrt2+t&s/2\\s/2&\sqrt2-t\end{pmatrix}
             \oplus(1-t)\oplus(1+t)\succeq0.
 \label{eq:tiltedlmi}
\end{equation}
The two-by-two determinant gives the ellipse, and the scalar blocks
give the interval for $t$.

The quantum upper boundary follows from the tilted-CHSH identity of
Bamps and Pironio~\cite{BampsPironio}. If
$T_\alpha=\alpha A_0+A_0(B_0+B_1)+A_1(B_0-B_1)$ and
$R_\alpha=\alpha A_1-A_0(B_0-B_1)-A_1(B_0+B_1)$, then
\begin{equation}
 T_\alpha^2+R_\alpha^2=(8+2\alpha^2)I,
 \qquad
 \beta I-T_\alpha=
 \frac{(\beta I-T_\alpha)^2+R_\alpha^2}{2\beta},
 \quad\beta=\sqrt{8+2\alpha^2}.
 \label{eq:tiltedsos}
\end{equation}
For a given $|t|\le1$, choosing $\alpha=2t/\sqrt{2-t^2}$ gives
$|s|\le2\sqrt{2-t^2}$.
Conversely, write $r=\sqrt{1-t^2}$ and $h=\sqrt{2-t^2}$.
The state and measurements
\begin{equation}
 |\psi_t\rangle=\sqrt{\frac{1+t}{2}}|00\rangle+
                \sqrt{\frac{1-t}{2}}|11\rangle,
 \quad A_0=Z,\quad A_1=X,\quad B_{0,1}=(Z\pm rX)/h
 \label{eq:tiltedattain}
\end{equation}
give $(t,s)=(t,2h)$. Reversing both Bob outcomes gives $(t,-2h)$;
mixing these experiments fills the vertical interval. This proves
both inclusions in Eq.~\eqref{eq:tiltedset}, including $t=\pm1$.
Theorem~\ref{thm:main} extends finite semidefinite representability
to all eight coordinates at once.

\section{Discussion}

The complete quantum set for two binary measurements per party has a
finite semidefinite representation with real algebraic coefficients.
The result places all Bell objectives in this scenario on one exact
finite optimization domain. It includes biased outcomes, mixed states
and arbitrary finite-dimensional implementations.

The proof depends on retaining the singular limits of the measurement
geometry. The real-qubit reduction makes the set semialgebraic; the
separated sum-of-squares construction supplies the additional positivity
that a finite lift requires. Independent coefficient and measurement
parameters expose the cases in which ordinary contact curvature
degenerates. Exact rescalings recover those directions while preserving
the complete state polynomial.

Finding a small explicit pencil is the next concrete problem.
The algebraic-coefficient conclusion makes exact coefficient search
possible in principle, while the local identities identify the
functions from which a constructive representation may be built.
The existence theorem itself is independent of the choice of a
moment hierarchy. Its conclusion is that the full quantum boundary
in the simplest Bell experiment admits one finite semidefinite
description.

\begin{thebibliography}{9}
\bibitem{NPA} M. Navascu\'es, S. Pironio, and A. Ac\'in,
A convergent hierarchy of semidefinite programs characterizing the set of quantum correlations,
New J. Phys. \textbf{10}, 073013 (2008),
\href{https://arxiv.org/abs/0803.4290}{arXiv:0803.4290}.
\bibitem{Masanes} Ll. Masanes,
Extremal quantum correlations for $N$ parties with two dichotomic observables per site,
\href{https://arxiv.org/abs/quant-ph/0512100}{arXiv:quant-ph/0512100} (2005).
\bibitem{BarizienBancal} V. Barizien and J.-D. Bancal,
Quantum statistics in the minimal scenario,
Nature Phys. \textbf{21}, 577 (2025),
\href{https://arxiv.org/abs/2406.09350}{arXiv:2406.09350}.
\bibitem{Shadows} C. Scheiderer,
Spectrahedral shadows,
SIAM J. Appl. Algebra Geom. \textbf{2}, 26--44 (2018),
\href{https://arxiv.org/abs/1612.07048}{arXiv:1612.07048}.
\bibitem{Scheiderer} C. Scheiderer,
Second-order cone representation for convex sets in the plane,
SIAM J. Appl. Algebra Geom. \textbf{5}, 114--139 (2021),
\href{https://arxiv.org/abs/2004.04196v2}{arXiv:2004.04196v2}, Corollary~3.19.
\bibitem{Basu} S. Basu,
Algorithms in real algebraic geometry: A survey,
\href{https://arxiv.org/abs/1409.1534}{arXiv:1409.1534} (2014).
\bibitem{BampsPironio} C. Bamps and S. Pironio,
Sum-of-squares decompositions for a family of CHSH-like inequalities and their application to self-testing,
Phys. Rev. A \textbf{91}, 052111 (2015),
\href{https://arxiv.org/abs/1504.06960}{arXiv:1504.06960}.
\end{thebibliography}

\end{document}

